% Exact Q-model endomorphism action lemma, proposed by root and audited here.
\section*{The actual endomorphism action on an absolutely simple CM variety over Q}

Let $A/\mathbf Q$ be absolutely simple of dimension $g$, and
suppose $E=\operatorname{End}^0_{\overline{\mathbf Q}}(A)$
is a CM field of degree $2g$. Then $E/\mathbf Q$ is Galois.
Moreover the actual conjugation action on endomorphisms
\[
       a:G_{\mathbf Q}\longrightarrow
                         \operatorname{Aut}_{\mathbf Q}(E)
\]
is surjective and its kernel is $G_E$. In particular every
endomorphism is defined over $E$, and this is the exact
field defining all of them. At any completely split prime,
the actual Galois action is transitive on the $2g$
one-dimensional $E$-eigenspaces of $V_\ell A$.

\paragraph{The type and its right stabilizer.}
Let $L$ be the normal closure of $E$, put
$G=\operatorname{Gal}(L/\mathbf Q)$ and
$H=\operatorname{Gal}(L/E)$, and fix the inclusion
$E\hookrightarrow L\subset\overline{\mathbf Q}$.
The embeddings of $E$ are identified with the right
cosets $G/H$ by $gH\mapsto g|_E$. Let $S\subset G$
be the inverse image of the CM type occurring in
$\operatorname{Fil}^1H^1_{\mathrm{dR}}(A/\mathbf Q)
\otimes\overline{\mathbf Q}$. Thus $S$ is right
$H$-invariant and
\[
                           S\sqcup cS=G,                  \tag{CMQ.1}
\]
where $c$ is the central complex conjugation of the
Galois CM field $L$. Here this assertion about the normal
closure has an explicit field proof. The field $L$ is
the compositum of the conjugate CM fields $g(E)$.
Complex conjugation on each $g(E)$ is its intrinsic
CM involution, so every embedding carries it to the
CM involution on the image. Hence complex conjugation
commutes with $G$ on each generating subfield and
therefore on $L$. Its fixed field is totally real:
for $x$ fixed by $c$ and $g\in G$, one has
$c(gx)=g(cx)=gx$. Since $L$ contains the totally
imaginary field $E$, the involution is nontrivial.
Thus $L/L^{\langle c\rangle}$ is a totally
imaginary quadratic extension of a totally real field.

Its right stabilizer is exactly $H$. Indeed let
\[
                    K=\{k\in G:Sk=S\}.
\]
It contains $H$, and $c\notin K$ by (CMQ.1).
The fixed field $F=L^K$ is a CM subfield of $E$:
centrality of $c$ makes its restriction to $F$ a
nontrivial involution, its fixed field lies in the
totally real field $L^{\langle c\rangle}$, and no
embedding of $F$ is real because $c$ is central and
acts nontrivially. Right $K$-invariance says precisely
that the type of $E$ is the inverse image, under
$\operatorname{Hom}(E,\mathbf C)\to
 \operatorname{Hom}(F,\mathbf C)$, of a CM type of $F$.

If $K$ were strictly larger than $H$, this would give
an actual proper abelian subvariety of $A$. To see the
map explicitly, identify $H_1(A(\mathbf C),\mathbf Q)$
with the rank-one $E$-module $E$. The induced-type
condition says that its complex structure, expressed
as multiplication in $E\otimes\mathbf R$, belongs to
$F\otimes\mathbf R$. Choose an $F$-basis of $E$ and
project onto one coordinate. The resulting rational
idempotent $p:E\to E$ commutes with the complex
structure and has rational rank $[F:\mathbf Q]$.
Choose an integer $n>0$ clearing its denominators on
$H_1(A(\mathbf C),\mathbf Z)$. The map $np$ gives a
complex-linear homomorphism of the actual analytic
torus $A(\mathbf C)$. Its graph is an algebraic
subvariety of the projective product by Chow's theorem,
so it is an abelian-variety endomorphism. The base-change
map
$\operatorname{End}_{\overline{\mathbf Q}}(A)
 \to\operatorname{End}_{\mathbf C}(A_{\mathbf C})$
is an isomorphism: spread any complex endomorphism over
a finitely generated extension of $\overline{\mathbf Q}$,
which is regular, and apply the graph-descent assertion
of Milne's original Corollary~20.4(b), pp.~146--147.
This is also the rigidity argument obtained by spreading
over a connected base and using the discrete Hom functor.
The actual endomorphism $np$ therefore descends to
$\overline{\mathbf Q}$ and has image of dimension
$[F:\mathbf Q]/2$ strictly between zero and $g$.
Absolute simplicity excludes this. Hence
\[
                              K=H.                        \tag{CMQ.2}
\]

\paragraph{The semilinear de Rham equation.}
The actual action $a_\sigma$ on $E$ corresponds to
some $n_\sigma H\in N_G(H)/H$:
\[
                       a_\sigma=n_\sigma|_E.
\]
For an $E$-eigenvector $w$ of de Rham character
$\varphi$ in the rational Hodge filtration, applying
$\sigma$ coefficientwise to $e^*w=\varphi(e)w$ gives
\[
             a_\sigma(e)^*\sigma(w)
                    =\sigma(\varphi(e))\sigma(w).
\]
The new character is
$\sigma\circ\varphi\circ a_\sigma^{-1}$.
The Hodge filtration is defined over $\mathbf Q$,
so the corresponding type identity is
\[
                         \sigma S=S n_\sigma.             \tag{CMQ.3}
\]
The right action has the stated orientation because
composing $g|_E$ with $n|_E$ gives $(gn)|_E$.
This also verifies compatibility with the
homomorphism law for actual endomorphism conjugation.

If $\sigma$ fixes $L$ pointwise, the left side of
(CMQ.3) is $S$. Equations (CMQ.2)--(CMQ.3) force
$n_\sigma\in H$ and $a_\sigma=1$. Thus $a$ factors
through the finite group $G$. For this factored
homomorphism, its kernel is exactly
\[
                 H^*=\{g\in G:gS=S\},                    \tag{CMQ.4}
\]
the left stabilizer of $S$. In particular $H^*$ is
normal in $G$.

Normality turns left stabilization into right
stabilization. For $s\in S$ and $h\in H^*$, write
$sh=(shs^{-1})s$. Since $shs^{-1}\in H^*$, the
right side belongs to $S$. Hence $SH^*=S$ and
\[
                                H^*\subset H.              \tag{CMQ.5}
\]
Now the exact index comparison is
\[
 [G:H]\leq [G:H^*]=|\operatorname{im}a|
       \leq|N_G(H)/H|\leq[G:H].                            \tag{CMQ.6}
\]
Every inequality is an equality. Thus $H^*=H$,
$N_G(H)=G$, and the endomorphism action is onto
$N_G(H)/H$. The subgroup $H$ is normal, so $E/\mathbf Q$
is Galois. Since $L$ was its normal closure, $L=E$
and $H=1$. Formula (CMQ.4) therefore proves that the
kernel of the original action is exactly $G_E$.

\paragraph{The actual Tate-module permutation.}
At a completely split prime let
$V_\varphi\subset V_\ell A$ be the eigenline for
$\varphi:E\to\mathbf Q_\ell$. The arithmetic action
is $\mathbf Q_\ell$-linear, and its endomorphism
conjugation gives the exact formula
\[
 \rho_\ell(\sigma)V_\varphi
                    =V_{\varphi\circ a_\sigma^{-1}}.       \tag{CMQ.7}
\]
Surjectivity of $a$ makes the orbit all the
$2g$ embeddings. A polarization over $\mathbf Q$
has Rosati involution equal to complex conjugation
on $E$. Its alternating pairing is therefore zero
between all eigenlines except the pairs
$V_\varphi,V_{\bar\varphi}$, and is nondegenerate
on each such pair. Scaling their generators yields
an actual symplectic basis made of the orbit lines.
This supplies both the arithmetic permutation and
the pairing needed by the simultaneous normalizer
calculation.

\paragraph{Exact scope for a proper geometric CM factor.}
For a proper CM quotient $q:J_Y\to B$ obtained only
over a finite extension, the displayed proof applies
to $B$ only after an actual $\mathbf Q$-model is
established. The equality (CMQ.3) used the rational
de Rham filtration of that same abelian variety.
The actual quotient map and its Galois conjugates
instead give maps to the factors $B^\sigma$;
their images remain subspaces of the specified
$V_\ell J_Y$. No replacement of that module by
$V_\ell B$ as the entire curve homology is made.
